% Auto-generated by scripts/06_emit_structure.py
% Source: scans 188–193
% Section id: appendix/4
% Phase 3 deliverable. Footnotes (Phase 4), citations (Phase 5),
% and Wade-Giles → Pinyin (Phase 6) are not yet applied here.

\chapter{Absolute Chronology: A Brief Note}
\label{app:4}

% source: scan 188, printed 171
% intro-echo-dropped: Absolute Chronology:
% intro-echo-dropped: A Brief Note
% intro-echo-dropped: I.
Introduction\footnote[1]{For a preliminary study of this problem, see Keightley (\citeyear{Keightley1975The}) and (\citeyear{Keightley1975bDate}).}
\appendixsectionlabel{sec:appendices-app04:1}{1}
The absolute date of the oracle-bone inscriptions is, strictly speaking, a subject for
historical rather than historiographical enquiry. It may be possible to deduce the date of the
inscriptions by a study of their content; it is not possible to do so by a study of their form or nature.
Since, however, certain of the statistical judgments proposed in appendixes 2 and 3 depend upon
the absolute length, if not the absolute dates, of the reigns from \pinyinterm{wu-ding} to \pinyinterm{di-xin}, a brief synopsis
of the issues and possible solutions will be given here.¹
2.
The Date of the Western Zhou Conquest
\appendixsectionlabel{sec:appendices-app04:2}{2}
Most attempts to determine the date of the Shang dynasty use the date of its fall as their
foundation. At least twenty different dates have now been proposed for the conquest, ranging
from 1122 to 1018 B.C., eloquent commentary on the inadequacies and ambiguities of the evidence
we possess.\footnote[2]{Barnard (\citeyear{Barnard1975First}), pp. 16-17, has the latest list (to which may be added Lao Gan [\citeyear{Lao1974Chou}] with a proposed conquest date of 1025 B.C.); cf. Barnard (\citeyear{Barnard1972Early}), pp. xxxix-xlii. Useful tabulations and discussions of the major theories may be found in Tung (\citeyear{Tung1945Yin}), pt. 1, ch. 2, pp. 6b-8a; ch. 4, pp. 11b-28b; (\citeyear{Tung1951bWuWang}), p. 199; Shima (\citeyear{Shima1966Bokuji}); Shirakawa (\citeyear{Shirakawa1971Kimbun}), pp. 280–281. As Chang Kwang-chih has remarked ([\citeyear{Chang1965Relative}], pp. 505-506), ``Some of these [dates], such as 1122, 1111, 1050, 1027, and 1018, may have been reconstructed on firmer grounds than others, but there is no one, in my opinion, that stands out and can be accepted as the most reliable."} %2
 \pinyinposs{sima-qian} failure to date the conquest suggests the difficulties involved. The
disagreements are essentially between two rival chronologies: the ``long'' one, which gives a conquest date in or near 1122; and the ``short'' one, which gives a conquest date in or near 1028 B.C.\footnote[3]{On the rival chronologies in general, see Chavannes (\citeyear{Chavannes1895Memoires}), pp. cxc-cxcvi; Gardner (\citeyear{Gardner1961Chinese}), p. 26, n. 8.} %3
The long chronology was first proposed by \pinyinterm{liu-xin} (first century B.C. and A.D.) in his
\pinyinterm{shijing-ages}, ``Canon of the Ages,'' a work created, with the aid of a smattering of historical
and quasi-historical records,\footnote[4]{For studies of \pinyinposs{liu-xin} \pinyinterm{shijing-ages}, see de Saussure (\citeyear{Saussure1924Chronologie}), pp. 322-339; Noda and Yabuuchi (\citeyear{Noda1945Kanjo}), pp. 137-179 and passim; \pinyinterm{chen-mengjia} (\citeyear{Chen1955ShangYin}), pp. 67-72.} %4
 to vindicate the accuracy of his \pinyinterm{santongli}, ``The Calendrical
% source: scan 189, printed 172
Complex of the Triple Concordance System,''\footnote[5]{This is the translation given by Sivin (\citeyear{Sivin1969Cosmos}), p. 12.} %5
 which was compiled about the time of Christ.
The date of 1122 B.C. was calculated by \pinyinterm{liu-xin} to accord with records about the supposed (and,
in fact, erroneous) position of Jupiter at the time of the conquest.\footnote[6]{The passage in question reads, ``From the Superior Epoch of the Triple Concordance to the year when (Wu Wang) attacked Chou (\pinyinterm{di-xin}, K28), there were 142.109 years. Jupiter was in Ch'un-huo, 13 degrees in (the lunar lodge) Chang'' (\pinyinterm{hanshu}, ch. 21b, p. 53b). For comment on this passage, see de Saussure (\citeyear{Saussure1924Chronologie}), pp. 330-331; Noda and Yabuuchi (\citeyear{Noda1945Kanjo}), p. 289.} %6
 It is clear from the errors they
contain that the records about the planet's position, attributed to early Chou tradition, depend on
retrospective, astronomical calculations made in Warring States or Han times.\footnote[7]{Eberhard, Müller, and Henseling (\citeyear{Eberhard1970Astronomie}), pp. 949-979; Noda and Yabuuchi (\citeyear{Noda1945Kanjo}), pp. 161-164.} %7
 They have no
astronomical validity; nor do the chronological conclusions that \pinyinterm{liu-xin} drew from them. Equally
flawed are \pinyinposs{liu-xin} attempts to relate his faulty Jupiter cycle to a Western Zhou lunar calendar,
reconstructed on Eastern Zhou principles applied to ambiguous passages of dubious historicity in
\booktitle{\pinyinterm{shangshu}, ``\pinyinterm{wu-cheng}''} and its preface.\footnote[8]{\pinyinterm{hanshu}, ch. 21b, pp. 60a-61a. The text has been glossed by de Saussure (\citeyear{Saussure1924Chronologie}), pp. 330-331; Noda and Yabuuchi (\citeyear{Noda1945Kanjo}), p. 289.} %8
 The details must be reserved for Studies, but it may
safely be stated that no reliable Zhou or Han data support the 1122 B.C. date.
The short chronology, which places the conquest in 1028 B.C., is based upon an entry in \pinyinterm{guben-zhushu-jinian}: ``From Wu Wang's extinguishing Yin down to Yu Wang was a total of 257 years.''\footnote[9]{Fan (\citeyear{Fan1962Kuben}), p. 35.} %9
This figure is no more reliable than \pinyinposs{liu-xin}. In the first place, we cannot be sure that the figure
of 257 years was in the original Jinian rather than in a subsequent commentary. Second, we cannot
be sure that the figure has been transmitted accurately. Third, it is uncertain if the 257 years refers
to the end or the start of Yu Wang's reign; if to the start, then the conquest must be advanced
eleven years to 1039 B.C. Fourth, and most importantly, there is no reason to think that the Jinian
author(s) had access to a chronological tradition that derived from accurate Western Zhou records.\footnote[10]{The chronological problems are discussed in detail by Barnard (\citeyear{Barnard1960bReview}). His study is particularly valuable for noting that the chronological data given in the Jinian are inconsistent and that ``the figure 257 years is insufficient to cover the absolute minimum of [Western Zhou] reign lengths as incompletely recorded in the Ancient Bamboo Annals'' (p. 506). By using internal evidence and by comparing the text itself with actual documents from the period which it purports to record, it may be shown that the Jinian is not a reliable source for historians, but a late Zhou view of early history. Its dubious historicity will be discussed more fully in Studies (for a preliminary treatment, see Keightley [\citeyear{Keightley1975aThe}]).} %10
Unlike \pinyinposs{liu-xin} 1122 date, the Jinian date cannot be rejected as spuriously determined, for we
have no idea how it was determined.
Some of the modern chronologies proposed for the Shang historical period are presented in
table 37. To a greater or lesser extent they all rely on the ``long'' or ``short'' chronologies and are,
in my view, flawed accordingly.\footnote[11]{Barnard (\citeyear{Barnard1975First}), pp. 30-31.} %11
3. Scientific Dating
\appendixsectionlabel{sec:appendices-app04:3}{3}
Barnard has argued that carbon-14 dates, dendrochronologically corrected, of \pinyinterm{anyang}
material (table 2) indicate ``fairly decisively'' the superiority of the long chronology over the short,
% source: scan 190, printed 173
but such a conclusion is premature for a number of reasons: \listitem{1} only two samples of charcoal
have been so dated;\footnote[12]{The Chinese reports give no information about how the charcoal was so identified. On the possibilities for, and consequences of, mistaken identification, see Michels (\citeyear{Michels1973Dating}), pp. 158-159.} %12
 \listitem{2} the presence of the charcoal (which might have lain on the ground for a
considerable period prior to burial) is unexplained; \listitem{3} the relative date of the sites-one, the large
tomb at Wu-kuan-ts'un,\footnote[13]{Sample ZK-5 was found in the great tomb, WKGM 1, excavated in 1950. The tomb's relative position within the historical period has not been established. The original report (\pinyinterm{guo-baojun} [\citeyear{KuoBaojun1951Yichiu}]) makes no attempt to do so. Barnard (\citeyear{Barnard1975First}), p. 31, refers to it as ``a comparatively late burial,'' but Kane (\citeyear{Kane1975Reexamination}), p. 109, tentatively dates it to ``\pinyintext{anyang} third generation,'' that is, by her count (which includes \pinyinposs{pangeng} generation), to Tung's period II. If we suppose, for sake of argument, that it was a period II burial, we still do not know if the charcoal was contemporary with the grave itself, or older. For the carbon-14 datings proposed for this sample, see table 2.} %13
 the other a cache of pyromantic scapulas\footnote[14]{The oracle bones with which sample ZK-86 was found are difficult to periodize with certainty. Kuo Mo-jo (\citeyear{Kuo1972Anyang}), p. 5, places them in period I; my own study of the ancestral titles suggests that they belong to the RFG, which I would also date to period I (see ch. 2, n. 18). Ch'iu (\citeyear{Chiu1972DuAnyang}) argues, unconvincingly in my view, for a period III + IV date. Whatever the date of the oracle-bone inscriptions (which could not, in any event, yield more than a terminus ante quem for the find), the date of the ceramic vessels and fragments from the same stratum must also be considered. At present, we know only that the excavators dated these to ``the first half of the late Yin period,'' which Kuo Mo-jo interprets to mean period III or IVa. Even if this dating is accepted and the pit is taken as a period III + IV find, it is still not clear if the charcoal should be dated to the time the scapulas were divined (period I?) or the time they were laid in the ground (period III + IV?, or neither).} %14
-is uncertain: it is
hard, therefore, to relate the carbon-14 dates to Tung's five periods; \listitem{4} the standard deviation is
so large (over ± 150 years, after dendrochronological correction, with even a one-sigma confidence
interval) that the two dates are useless for helping us choose between the long and short chronologies;\footnote[15]{The statistical issues are explained by Clark (\citeyear{Clark1975Calibration}), pp. 265-266. If, as he advises, one works with 95 percent (or two-sigma) confidence intervals, then the intervals that result are too gross to be useful for the judgments we need to make (see table 2). But if, for two samples---such as ZK-5 and ZK-86---one uses 68 percent (or one-sigma) confidence intervals, as Barnard has done, then there is roughly a fifty-fifty chance that ``the actual calendar date of at least one of the samples will lie outside the corresponding confidence intervals.''} %15
 \listitem{5} the half-life of carbon-14 is still not determined beyond doubt;\footnote[16]{E.g., Goodyear (\citeyear{Goodyear1971Archaeological}), p. 181: ``More serious doubts surround the precise value of the half-life of C4, which different workers have put at from 4,700 to 7,200 years, but which has now been provisionally accepted internationally as 5,568 ± 30 years. It is not possible to give any one definite figure for the reliability of radiocarbon dates, but one must allow for a possible error of at least 10 percent.'' Clark (\citeyear{Clark1975Calibration}) uses the 5568 half-life. Barnard (\citeyear{Barnard1975First}) usually records both the 5568 and the 5730 half-life figures.} %16
 \listitem{6} the degree to
which dendrochronological corrections, based largely on data from other parts of the world, may
safely be applied to the Chinese Bronze Age is not yet certain;\footnote[17]{Cf. Barnard (\citeyear{Barnard1975First}), p. 38. To the best of my knowledge, the only Far Eastern data incorporated into discussions of the accuracy of dendrochronological corrections are those of Kigoshi and Hasegawa (\citeyear{Kigoshi1966Secular}), who studied tree rings from a giant cedar taken from Yaku Island, Japan; Clark (\citeyear{Clark1975Calibration}), p. 260, makes use of their work.} %17
 \listitem{7} the analytical and statistical
techniques of the Chinese carbon-14 laboratories have not yet been compared with those elsewhere;
the question of interlab variability, therefore, introduces another unknown factor.\footnote[18]{Cf. Barnard (\citeyear{Barnard1975First}), pp. vi, vii, 21-22; Clark (\citeyear{Clark1975Calibration}), pp. 252, 258-259.} %18
 We will need
scores of carbon-14 dates from the \pinyinterm{anyang} site before any firm conclusions based on this method
can be drawn about the absolute dates of the historical period. At present, any conclusions based
on the limited evidence we have rest upon a large number of uncertain assumptions.
% source: scan 191, printed 174
4. Celestial Phenomena
\appendixsectionlabel{sec:appendices-app04:4}{4}
The attempt to provide absolute dates for the lunar eclipses recorded in the oracle
bones also rests upon a number of uncertain assumptions. It assumes in particular: \listitem{1} that the
inscriptions are true astronomical records, not fabricated portents; and \listitem{2} that the sixty-day
kan-chih cycle, which the Shang used to date the days and record the eclipses, was in continuous
use, without error, from the reign of \pinyinterm{wu-ding} down to the present (this must be so if we are to
relate the free-floating cyclical dates to absolute dates). The matter will be discussed fully in
Studies, but there are, I believe, reasons for accepting these assumptions as true. Using a ``Canon
of Lunar Eclipses from - 1500 to -1000 with Conditions for Determining Visibility at Anyang''
developed by Robert R. Newton, four period I lunar eclipse records, all verifications, may be dated
to the eighteen-year period from 1198 to 1180 B.C.\footnote[19]{This canon, which has now been published as (Newton [\citeyear{Newton1977Canon}]), supersedes that of Dubs (\citeyear{Dubs1947Canon}); Newton's canon is both more accurate, being based upon more recent studies of the movement of the earth, sun, and moon, and more realistic in that it enables us to allow for possible eccentricities in the orbits of the earth and moon which would lead to errors in the retrospective calculations. For a preliminary treatment of the eclipse inscriptions, see \textcite[pp.~132--174]{Keightley1975The}; \textcite{Keightley1975bDate}; for a recent study, which arrives at different conclusions, see Zhang Peiyu (\citeyear{Chang1975Jiaguwen}). The four verifications and the dates I would assign them are: \pinyinterm{pinbian} 57.1 (4 November 1198); \pinyinterm{fuyin}, ``T'ien'' 2 (25 October 1189); Jinzhang 594 (D) and \pinyinterm{kufang} 1595 (D) (27 December 1192; note that the inscription itself can be dated to the Shang twelfth month); and \pinyinterm{pinbian} 59.1/60 (22 May 1180); these inscriptions are transcribed at S162.1. A fifth lunar eclipse may have been recorded as a divination charge on one of the scapulas (H23:66) discovered in 1973. The inscription reads, in part, \pinyinterm{renyin} *tieng yüeh yu chih (*fiak). If *fiak was equivalent to *diak---\pinyinterm{houbian} 1.29.6 and \pinyinterm{cuibian} 55 [S581.2] offer good evidence that this was so; cf. Shima [\citeyear{Shima1958Inkyo}], p. 270; Serruys [\citeyear{Serruys1974The}], p. 104---we may translate: ``On \pinyinterm{renyin} (day 39) divining: `The moon was eclipsed.' '' If we assume that the (daytime?) charge referred to an eclipse in the night before (day 38) and if, with the authors of the preliminary report (KK [\citeyear{KK1975Anyang}], p. 38; cf. KK [\citeyear{KK1975aJufa}], p. 142 and plate 1), we date the inscription to period IVa, then this could be a record of the eclipse of 28 December 1127 B.C., providing added support for the revised short chronology of table 38, which dates period IVa to about 1130-1116 B.C. Uncertainties about the meaning of the inscription, about the cyclical date of the eclipse, and about the period of the inscription, however, make this record less reliable than the four cases cited above.} %19
 Other, less plausible, clusters of dates may be
assigned, but if the eclipses are dated correctly-and the fact that the first and last of the four
period I eclipses were recorded on plastrons from the same pit, YH127, whose contents would
thus have spanned some eighteen years, offers some support for this view-then the short chronology is to be preferred to the long (cf. sec.~\ref{sec:appendices-app04:5} below% was: (cf. sec. 5 below)
).
5. Average Reign Lengths
\appendixsectionlabel{sec:appendices-app04:5}{5}
Only in period V does the combination of a rigid ritual schedule with the recording of
numbered ritual cycles (sec.~\ref{sec:chapters-ch04:4.3.1.7}% was: (sec. 4.3.1.7)
) permit educated estimates to be made of the reign lengths
of \pinyinterm{di-yi} and \pinyinterm{di-xin}. These estimates vary in accordance with the way particular scholars have
reconstructed the period V ritual calendar. It is clear from the inscriptions, however, that the
reigns of these two kings each lasted from twenty to thirty-three years.\footnote[20]{\pinyinterm{chen-mengjia} (\citeyear{Chen1955ShangYin}), p. 59. Shima (\citeyear{Shima1960TiYi}), p. 51, argues that \pinyinterm{di-yi} was succeeded by \pinyinterm{di-xin} in the seventh month of \pinyinposs{di-yi} twenty-first cycle; subsequently ([\citeyear{Shima1966Bokuji}], p. 13), he located the reign change in \pinyinposs{di-yi} twentieth cycle. \pinyinterm{xu-jinxiong-tight} (\citeyear{Hsu1970Wuzhong}), p. 26, using a new interpretation of the ritual cycle derived from inscriptions not used by Shima, argues that \pinyinterm{di-yi} reigned 31 years. Tung (\citeyear{Tung1945Yin}), pt. 1, ch. 3, p. 18b; pt. 2, ch. 1, p. 42b, allotted 35 years to \pinyinterm{di-yi}. As to \pinyinterm{di-xin}, Shima (\citeyear{Shima1960TiYi}), p. 52, first argued that he was on the throne for 30-plus cycles, then refined that figure ([\citeyear{Shima1966Bokuji}], p. 21; cf. [\citeyear{Shima1976Teishin}]) to argue that \pinyinterm{di-xin} was killed during the thirty-third cycle, and thus reigned 32 cycles. Tung (\citeyear{Tung1945Yin}), pt. 1, ch. 3, p. 19a; pt. 2, ch. 1, p. 50b, by contrast, allotted 63 years to \pinyinterm{di-xin}. Since no inscriptions which record more than a twentieth cycle have been found, this figure is implausibly large.} %20
 For reign lengths in
periods I to IV, the inscriptions contain no information of value.
% source: scan 192, printed 175
It is possible, however, to estimate the length of the average reign. On the basis of the reign
lengths recorded in \pinyinterm{shiji}, Bishop arrived at an average reign length of twenty years for the
various kings and state rulers of Chou.\footnote[21]{Bishop (\citeyear{Bishop1932Chronology}), pp. 234-235. Similar results may be obtained from the approach of Lei (\citeyear{Lei1931Yinzhou}), who calculated that the average royal generation throughout Chinese history lasted twenty-five years. If we assume that the lunar eclipse inscriptions place \pinyinterm{wu-ding} on the throne at about 1200-1180 B.C., and that there were seven generations of Shang kings during the historical period, we arrive at a conquest date of (1200 - [7 x 25]) = ca. 1025 B.C. Lei, working backward from known Chou dates, arrived at a conquest date of about 1030-1020 B.C. (cf. Ho [\citeyear{Ho1975Cradle}], p. 4). Whether reign lengths are more accurate for our purposes than generations is problematical.} %21
 It seems likely that the reigns of the historical Shang
period would have been of the same average length (if not shorter).\footnote[22]{It is unlikely that the Shang kings were markedly longer-lived than the kings of Western Zhou. This consideration alone challenges many of the traditional reign lengths given in table 37; cf. n. 24 below. \pinyinterm{guben-zhushu-jinian} (Fan [\citeyear{Fan1962Kuben}], p. 24) tells us that from the time T'ang (K1) extinguished the Xia down to Shou (K28), there were twenty-nine Shang kings in 496 years. For what this is worth---and all the objections raised at n. 10 above apply---this would give an average reign of seventeen years.} %22
 This suggests that the \pinyinterm{wu-ding}-\pinyinterm{di-xin} era would have lasted 160 years (if there were eight kings) or 180 years (if there
were nine).\footnote[23]{On the question of whether \pinyinterm{lin-xin} reigned or not, see table 1, note h.} %23
If we use Bishop's figure of twenty years per reign for the ten Western Zhou kings who ruled
prior to the start of the Kung Ho period in 841 B.C. (the first firm date we are thought to
possess), we can estimate that the conquest would have taken place about 1041 B.C.\footnote[24]{Cf. Bishop (\citeyear{Bishop1932Chronology}), pp. 232-237. The long chronology, which places the conquest in 1122, would require that the ten kings have ruled for an average of twenty-eight years (cf. n. 22). The average reign lengths of all but one Chinese dynasty have been far shorter; for example, Eastern Zhou is estimated at 22 years; Sung, at 21 years; Later Han, at 16 years; T'ang, at 12 years, etc. It is unlikely that the Western Zhou kings would have reigned, on the average, nearly as long as the Ch'ing emperors (about thirty years). It is not entirely certain, of course, that the figure of ``ten kings'' is accurate, though all reconstructions of Western Zhou chronology (as listed by Barnard [\citeyear{Barnard1975First}], pp. 16-17) use it.} %24
 This figure
accords well with the dates provided by the eclipse inscriptions, which place \pinyinterm{wu-ding} on the
throne from about 1200 to 1180 B.C. If (deleting \pinyinterm{lin-xin}),\footnote[25]{See n. 23.} %25
 seven more kings ruled before the
conquest, and if those seven kings reigned for an average of twenty years, we arrive at a conquest
date of 1040 B.C.
6.
Conclusions: A Revised Short Chronology
\appendixsectionlabel{sec:appendices-app04:6}{6}
The revised short chronology presented in table 38 is a tentative one, based on the lunar
eclipse dates and average reign lengths, the latter adjusted to take account of information found
in the inscriptions. It is offered only as a working model, and it is undoubtedly inaccurate in its
theoretical rationing of the reign lengths. It has the advantage of being based on no legendary
information such as the fifty-nine years traditionally assigned to \pinyinterm{wu-ding}\footnote[26]{See table 37, note d.} %26
---yet it does accord
% source: scan 193, printed 176
well with many Eastern Zhou and Han traditions.\footnote[27]{As \pinyinterm{chen-mengjia} ([\citeyear{Chen1955ShangYin}], pp. 55-56) has shown, passages in \pinyinterm{zuozhuan}, \pinyinterm{mengzi}, and \pinyinterm{shiji}, ``\pinyinterm{lu-shijia},'' suggest that the authors of these texts believed the Western Zhou to have lasted for about 240 to 290 years (which would give a conquest date of 1060-1010), with a length of 250 to 260 years (and thus a conquest date of 1030-1020 B.C.) being favored. The reliability and usefulness of these figures is uncertain-on the \pinyintext{mengzi} passage in particular see Knoblock (1964).} %27
 I doubt that any of its figures are in error
by more than twenty-five years. Given the disproportionately large number of period I inscriptions
in the corpus (sec.~\ref{sec:chapters-ch05:5.3.2}% was: (sec. 5.3.2)
), it is possible that \pinyinterm{wu-ding} reigned longer than twenty years; but whether
that would put the start of his reign significantly before 1200 B.C. or the end of his reign well after
1180 B.C. is problematic. It may be noted, incidentally, that the reign dates given in table 38 are
not excluded by the initial carbon-14 determinations given in table 2.\footnote[28]{ZK-5 is said to come from a period II burial (see n. 13). It should, according to table 2, using one-sigma confidence intervals, be carbon-14 dated to 1403-1079 B.C.; the revised short chronology of table 38 assigns a date of 1180-1151 B.C. to period II. According to table 2, ZK-86 should be carbon-dated to 1429-1124 B.C.; if this sample is, at the latest, a period III + IV charcoal fragment (see n. 14), then the revised short chronology (table 38) would assign a date of 1150-1091 B.C. to period III + IV. The possibility that the charcoal was older than the burial has already been mentioned.}
